% Copyright (c) 2022-2026 MKM Research Labs. Licensed under the MIT License; see LICENSE.
%
% MKM-ModelRisk — Integration Guide.
%
% Build:   tectonic docs/integration_guide.tex      (or: latexmk -pdf docs/integration_guide.tex)
%
% The three colours below are the values held in config/theme.py (BRAND accent,
% TEXT muted, RAG red). This file is written by hand rather than generated by the
% platform, so it is outside the R10 styling audit's scope (src/static,
% src/templates and .py under config/src/scripts/app) — but it follows the same
% discipline: if the adopter rebrands, these three lines change with the theme.
\documentclass[11pt]{article}

\usepackage[utf8]{inputenc}
\usepackage[T1]{fontenc}
\usepackage{geometry}\geometry{a4paper, left=22mm, right=22mm, top=22mm, bottom=25mm}
\usepackage{graphicx}\usepackage{longtable}\usepackage{array}\usepackage{booktabs}
\usepackage{xcolor}\usepackage{parskip}\usepackage{enumitem}\setlist{nosep}
\usepackage{float}\usepackage{lastpage}\usepackage{fancyhdr}
\usepackage{amsmath}\usepackage{amssymb}
\usepackage{fancyvrb}
\usepackage{tcolorbox}
\usepackage{hyperref}
% NOT hyphenat: it breaks long \texttt{} names with a hyphen, and a hyphen
% inside DB_APP_USER or mkm_modelrisk_scratch is a name a reader might copy.
% Long identifiers are broken by hand where a column cannot hold them.
\usepackage{microtype}
\setlength{\emergencystretch}{3em}
\setcounter{tocdepth}{2}\setcounter{secnumdepth}{3}

% --- MKM Research Labs house style (mirrors src/governance/latex.py) ---
% When the console publishes this guide into the Audit Reports tab it writes
% theme_colours.tex beside the source, from config/theme.py via config/reporting.py,
% so the published copy carries the adopter's brand (rule R10). A build by hand has
% no such file and takes the shipped values below; \providecolor leaves whatever the
% input already defined alone.
\IfFileExists{theme_colours.tex}{\input{theme_colours.tex}}{}
\providecolor{mkmblue}{HTML}{1976D2}
\providecolor{mkmchrome}{HTML}{888888}
\providecolor{mkmred}{HTML}{D32F2F}
\providecolor{mkmwash}{HTML}{F7FBFF}
\providecolor{mkmline}{HTML}{E0E0E0}
\hypersetup{colorlinks=true,linkcolor=mkmblue,citecolor=mkmblue,urlcolor=mkmblue,
            pdftitle={MKM-ModelRisk Integration Guide},pdfauthor={MKM Research Labs}}

\pagestyle{fancy}\fancyhf{}
\fancyhead[L]{\small\textcolor{mkmchrome}{MKM Research Labs}}
\fancyhead[R]{\small\textcolor{mkmchrome}{MKM-ModelRisk --- Integration Guide}}
\fancyfoot[C]{\small\textcolor{mkmchrome}{CONFIDENTIAL --- PROPRIETARY}}
\fancyfoot[R]{\small\textcolor{mkmchrome}{Page \thepage\ of \pageref{LastPage}}}
\renewcommand{\headrulewidth}{0.4pt}\renewcommand{\footrulewidth}{0.4pt}

% Verbatim blocks, set small with a rule down the left so a listing is visibly
% not prose.
\DefineVerbatimEnvironment{code}{Verbatim}%
  {fontsize=\small, frame=leftline, framerule=0.8pt, rulecolor=\color{mkmline},
   framesep=6pt, xleftmargin=6pt, baselinestretch=1.0}

\newtcolorbox{note}[1][]{colback=mkmwash, colframe=mkmblue, boxrule=0.6pt,
  left=6pt, right=6pt, top=5pt, bottom=5pt, arc=2pt, title={#1},
  fonttitle=\bfseries\small, coltitle=white}

\newcommand{\cfg}[1]{\texttt{#1}}
\newcommand{\rulename}[1]{\textbf{#1}}

\begin{document}

% --- Title page (mirrors latex.py title_page + legal_notice) ---
\begin{titlepage}\begin{center}
\vspace*{2cm}
{\Huge\bfseries\textcolor{mkmblue}{MKM-ModelRisk}\par}\vspace{0.6cm}
{\LARGE\textcolor{mkmchrome}{Integration Guide}\par}\vspace{1.2cm}
{\large How the platform works, how to restyle it, how the coding rules\par
 are enforced, how to move it to another SQL engine or a container,\par
 and how it takes identity from your SSO.\par}\vspace{2.5cm}
{\large MKM Research Labs\par}\vspace{0.4cm}
{\large Version 1.0 --- \today\par}\vspace{0.4cm}
{\large Written against \texttt{main} at commit \texttt{642d4f6} (2026-08-23)\par}
\vspace{2cm}\textcolor{mkmred}{\textbf{CONFIDENTIAL --- PROPRIETARY}}
\end{center}
\vfill
\begin{minipage}{\textwidth}\footnotesize\raggedright
\textbf{Legal notice}\par\medskip
PROPRIETARY AND CONFIDENTIAL\par\medskip
This document describes the MKM-ModelRisk platform. The software itself is
MIT-licensed (see \texttt{LICENSE} in the repository); this notice covers the contents of
this document, which constitute the confidential material of MKM Research Labs
(``MKM'').\par\medskip
Any usage, reproduction, distribution, modification, or disclosure of this document or any
portion thereof, without the express prior written authorisation from MKM Research Labs is
strictly prohibited.\par\medskip
The document is provided on an ``as is'' basis. MKM Research Labs makes no representations
or warranties of any kind, express or implied, regarding its accuracy, reliability, or
suitability for any particular purpose.\par\medskip
All rights reserved. \textcopyright{} 2022--2026 MKM Research Labs.\par\medskip
Governed by the laws of the United Kingdom.
\end{minipage}
\end{titlepage}

\tableofcontents
\clearpage

\section{About this guide}

\subsection{Who it is for}

This guide is for the engineer who has to \emph{take MKM-ModelRisk into an institution}:
stand it up, point it at that institution's database, put its own brand on it, wire it to
its single sign-on, and satisfy whoever asks how the thing is controlled. It assumes you
can read Python and SQL. It does not assume you know model-risk governance --- for what
the screens \emph{mean}, read \texttt{docs/user\_guide.md}; this guide is about the
machinery underneath them.

\subsection{What it covers}

\begin{enumerate}[label=\S\arabic*.]
  \setcounter{enumi}{1}
  \item How the technology works --- the event store, the read models, the request path,
        and where configuration comes from.
  \item Changing the style --- the single file that holds every colour, size and typeface,
        how it reaches the browser and the PDFs, and how to rebrand safely.
  \item The coding rules R1--R10 --- what each one is for, which are mechanically gated,
        and how to run the audit.
  \item Switching to a different SQL server --- what is engine-specific, what is not, and
        the checklist for a port.
  \item Running it inside a container --- what has to be in the image, a worked
        \texttt{Dockerfile} and Compose file, and the traps that only appear once the app
        is not on a developer's laptop.
  \item Permissions and SSO --- who is allowed to author a change, where the trust
        boundary sits, and how to put an identity provider in front of it.
\end{enumerate}

\subsection{How to read the claims in it}

Every statement here is written against the code at the commit on the title page, and the
guide is explicit about the difference between three things that are easy to conflate:

\begin{itemize}
  \item \textbf{Built and enforced} --- there is code, and something fails if you break it.
  \item \textbf{Built, not enforced} --- there is code, and nothing stops you bypassing it.
  \item \textbf{Designed, not built} --- there is a design document and no code.
\end{itemize}

Where a section describes work you would have to do yourself (the SQL Server port, the
container image), it says so plainly rather than presenting a plan as a feature.

\clearpage
\section{How the technology works}

\subsection{One paragraph}

MKM-ModelRisk is a model-risk-governance platform: it holds an institution's inventory of
quantitative models, the validation and committee process around them, and the evidence
trail that process produces. Its distinguishing property is that \textbf{the audit trail is
the system of record}. Nothing is stored as a mutable row that somebody edited; every
change is an \emph{event} appended to a hash-chained log, and every screen you see is a
projection rebuilt from that log. The application is a Flask app serving a dependency-free
vanilla-JavaScript console, on PostgreSQL, with PDFs rendered through \LaTeX.

\subsection{The shape of a request}

\begin{code}
  browser (vanilla JS console)
        |  fetch /api/...  + X-Auth-Request-User header
        v
  Flask blueprint  (src/web/api_*.py)
        |  context_from_request(): resolve_principal -> Principal
        v
  dispatch()       (src/command/dispatch.py)
        |  1. confirm-to-amend gate   2. authorize()   3. actor + reason
        v
  domain command   (src/governance/**)
        |  append_event(...)                    read_view(...)
        v                                            ^
  src/database/    ---- events (append-only) --------+
        |                     |
        |                     +--> projections update the *_view read models
        v
  PostgreSQL
\end{code}

Four things about that path are worth understanding before you change anything.

\paragraph{Every mutation goes through one seam.} \texttt{dispatch()} is the only sanctioned
way to run a command that changes state. It enforces confirm-to-amend (the UI's
confirmation dialog cannot be bypassed by calling the API directly), calls
\texttt{authorize()}, and threads the resolved actor and the change reason into the
handler, which stamps them on the event. A new endpoint that forgets to ask about
permissions is not possible so long as it dispatches; the few paths that cannot dispatch
--- multipart uploads, which call their command directly --- apply the same check through
\texttt{request\_actor()} in \texttt{src/web/\_helpers.py}.

\paragraph{Every persistence call goes through one package.} No module outside
\texttt{src/database/} may contain SQL, a cursor, a connection or a driver import. That is
coding rule \rulename{R6}, it is audited with zero tolerance, and it is the single reason
\S\ref{sec:sqlserver} (moving to another SQL engine) is a bounded job rather than an
open-ended one.

\paragraph{The log is append-only, and the database says so.} Migration
\texttt{0002\_roles\_grants.sql} grants the runtime role \texttt{SELECT} and
\texttt{INSERT} on \texttt{events} and revokes \texttt{UPDATE}, \texttt{DELETE} and
\texttt{TRUNCATE}. Immutability is therefore a property of the deployment's grants, not a
convention in the application. Each event also carries \texttt{prev\_hash} and
\texttt{event\_hash} --- a SHA-256 chain --- so tampering by a more privileged account is
\emph{detectable} even where it is not \emph{prevented}.

\paragraph{Read models are caches, and are rebuilt on boot.} The \texttt{*\_view} tables
are derived from the log and maintained incrementally. If one drifts --- after a restore,
or a data-directory move, where no command ran to refresh it --- records render blank even
though the log is intact. This bit in production on 2026-07-29: an MRC meeting looked
deleted when only its view had gone stale. Startup therefore replays every projection, so
a stale cache cannot outlive a restart. Set
\cfg{REBUILD\_VIEWS\_ON\_STARTUP=0} to skip it if the estate ever grows large enough that a
full replay per boot is undesirable.

\subsection{The event store}

The schema is deliberately small. \texttt{src/database/migrations/0001\_event\_store.sql}:

\begin{code}
CREATE TABLE IF NOT EXISTS events (
    global_seq     BIGINT GENERATED ALWAYS AS IDENTITY PRIMARY KEY,
    event_id       UUID NOT NULL DEFAULT gen_random_uuid(),
    aggregate_type TEXT NOT NULL,
    aggregate_id   TEXT NOT NULL,
    stream_seq     INTEGER NOT NULL CHECK (stream_seq > 0),
    event_type     TEXT NOT NULL,
    payload        JSONB NOT NULL DEFAULT '{}'::jsonb,
    actor          TEXT NOT NULL,
    reason         TEXT,
    recorded_at    TIMESTAMPTZ NOT NULL,
    prev_hash      TEXT,
    event_hash     TEXT NOT NULL,
    CONSTRAINT events_stream_unique UNIQUE (aggregate_type, aggregate_id, stream_seq)
);
\end{code}

A \emph{stream} is one governed subject --- a model, a product, a committee meeting ---
keyed by \texttt{(aggregate\_type, aggregate\_id)}. \texttt{stream\_seq} is a gap-free
1-based counter within it, and the unique constraint is the optimistic-concurrency
backstop: two writers cannot both claim position \emph{n}. \texttt{global\_seq} orders the
whole log, and each projection records how far it has replayed in
\texttt{projection\_checkpoints}.

The public API of \texttt{src/database} is intent-named --- \texttt{append\_event},
\texttt{read\_stream}, \texttt{read\_view}, \texttt{replay}, \texttt{run\_migrations} ---
and there is no escape hatch. If a query you need is not exposed, the correct move is to
add a named function to that package, not to open a connection where you are.

\subsection{Migrations}

Numbered plain-SQL files under \texttt{src/database/migrations/}, applied in lexical order
by the owner role and recorded in \texttt{schema\_migrations}, so re-running is a no-op.
Two details matter operationally:

\begin{itemize}
  \item \textbf{The runtime role is substituted, not hardcoded.} The grants are written
        against \texttt{\{\{app\_user\}\}} / \texttt{\{\{app\_password\}\}}, which the
        runner fills in from \cfg{DB\_APP\_USER} / \cfg{DB\_APP\_PASSWORD}. A deployment
        that renames the role still gets grants against the role it actually connects as.
        The name is validated against a plain-identifier pattern first, because a role name
        cannot be passed as a query parameter.
  \item \textbf{A PostgreSQL advisory lock is held while migrating.} Under gunicorn several
        processes may start at once; without the lock two of them read the same migration
        as pending, both execute its DDL, and the loser dies on the
        \texttt{schema\_migrations} primary key --- a worker that crashes on boot with
        nothing in the log to explain it.
\end{itemize}

\subsection{Entry points and startup}

\begin{longtable}{@{}p{0.32\textwidth}p{0.62\textwidth}@{}}
\toprule
\textbf{Command} & \textbf{What it is} \\
\midrule
\endhead
\texttt{python model.py} & Development server. Werkzeug, single process, auto-reloading.\\
\texttt{gunicorn wsgi:application} & The real thing. Settings come from
  \texttt{gunicorn.conf.py}, which is a thin wiring of \texttt{config/web.py}. \\
\texttt{scripts/start.sh} & One-command setup on Linux/macOS: toolchain checks, venv,
  dependencies, local cluster, then launch. Idempotent. \\
\texttt{scripts/dev\_app.sh} & The verification app on port 5016 against the scratch
  database. Refuses to start if the resolved database is a protected one. \\
\bottomrule
\end{longtable}

The one-time startup work --- migrations, seeding the admin reference lists, and the
read-model rebuild --- is \texttt{bootstrap()} in \texttt{src/web/\_startup.py}. Under
gunicorn it runs in the \textbf{master}, in the \texttt{on\_starting} hook, before any
worker is forked; the master sets \cfg{MODELRISK\_BOOTSTRAP\_DONE} in the environment and
the forked workers inherit it and skip the work. Under any other WSGI server, importing
\texttt{wsgi.py} does it. Either way it happens exactly once per boot, before the first
request.

The hourly database backup thread is started in the same hook, for the same reason: one
schedule, owned by the process that outlives worker recycling.

\subsection{Where configuration comes from}

Everything parameterised lives in the \texttt{config/} package --- that is coding rule
\rulename{R1} --- and each module reads its values from the environment with a shipped default:

\begin{longtable}{@{}p{0.30\textwidth}p{0.64\textwidth}@{}}
\toprule
\textbf{Module} & \textbf{Owns} \\
\midrule
\endhead
\texttt{config/database.py} & Host, port, database name, the owner and app roles,
  the protected-database list, the migration tokens and lock key. \\
\texttt{config/web.py} & Bind address, worker model, timeouts, logging. \\
\texttt{config/identity.py} & Which request headers carry the authenticated identity. \\
\texttt{config/permissions.py} & Who may author a change; the role vocabulary. \\
\texttt{config/theme.py} & Every colour, type size, spacing step, radius and shadow. \\
\texttt{config/reporting.py} & How a generated PDF is set, and its shipped wording. \\
\texttt{config/files.py} & Upload directory, allowed extensions, size limits. \\
\texttt{config/paths.py} & Repository-relative paths (migrations, contracts, static). \\
\texttt{config/audit.py} & Audit thresholds and scan scope. \\
\texttt{config/backup.py} & The backup schedule. \\
\texttt{config/governance/} & The governance vocabulary and field metadata. \\
\bottomrule
\end{longtable}

\begin{note}[The \texttt{.env} trap --- read this before running anything ad hoc]
\texttt{config/env.py} loads \texttt{.env} \emph{at import time}, so \emph{any} process
that imports \texttt{config} inherits that file's settings --- including
\cfg{DB\_NAME} pointing at the live database. A throwaway
\texttt{python -c "from src.web import create\_app"} is aimed at production by default. A
real environment variable always beats the file, so the fix is to set the target
explicitly and print what you actually resolved to before the first write:

\begin{code}
python -c "from config.database import DatabaseConfig; \
           print(DatabaseConfig.from_env().name)"
\end{code}

This is coding rule \rulename{R9}, and it exists because two ad-hoc debugging scripts wrote
junk models into the live, append-only governance log on 2026-08-06. Those records are
permanent; removing them from the active inventory took a committee.
\end{note}

\subsection{The console and the documents}

The front end is deliberately dependency-free: vanilla JavaScript modules under
\texttt{src/static/js/}, one stylesheet at \texttt{src/static/css/app.css}, and a single
Jinja template \texttt{src/templates/index.html}. There is no build step, no bundler and no
framework to keep current --- a decision that matters for a system whose useful life is
measured against a regulator's retention period rather than a framework's release
cadence. Rule \rulename{R7} keeps those assets out of \texttt{src/web/}, and forbids
embedding JavaScript or CSS inside a \texttt{.py} file at all.

PDFs --- the MRC meeting pack, the Model Risk Report, the coding-rule audit reports, the
governance completeness audit --- are generated as \LaTeX{} and compiled by
\textbf{tectonic}, which must be on \texttt{PATH}. Without it those endpoints return
\texttt{503} and everything else works. The \LaTeX{} source of a generated document can be
filed alongside its PDF as the auditable original; uploads are only ever stored and served
back, never compiled, so no user-supplied \LaTeX{} is executed on the server.

\clearpage
\section{Changing the style}
\label{sec:style}

\subsection{The rule this section is about}

An adopting institution has to be able to put its own brand on the platform without
touching the front end. That is coding rule \rulename{R10}, added on 2026-08-22, and it is
absolute: \textbf{no colour literal --- hex, \texttt{rgb()}, \texttt{rgba()},
\texttt{hsl()}, or the bare \texttt{RRGGBB} spelling \LaTeX{} uses --- may appear anywhere
except \texttt{config/theme.py}}. Not in the CSS, not in a template, not in a
\texttt{.js} file, not in a \texttt{.py} file, and not in the PDF preamble.

The rule exists because the alternative had already failed. Before it, the brand blue
existed as \texttt{\#1976d2} in the stylesheet and the JavaScript and as
\texttt{\#1565C0} in the PDF preamble, in two letter-cases: one decision, written down four
times, already disagreeing with itself. A rebrand was a find-and-replace across three
languages and it would have missed some.

\subsection{Where the values live}

\texttt{config/theme.py} --- one file, around 280 lines, organised in groups:

\begin{longtable}{@{}p{0.20\textwidth}p{0.74\textwidth}@{}}
\toprule
\textbf{Group} & \textbf{What it holds} \\
\midrule
\endhead
\texttt{BRAND} & The adopter's identity: \texttt{accent} (links, active tabs, focus,
  selection), its soft/border/ink/wash variants, and the header gradient. \emph{The group a
  rebrand almost always touches, and often the only one.} \\
\texttt{SURFACE} & Paper: background, panel, raised, sunken, control, read-only, code,
  and the two rule weights. \\
\texttt{TEXT} & Ink: primary, secondary, muted, inverse. \\
\texttt{RAG} & Red / Amber / Green. This is the \emph{product's vocabulary}, not a colour
  choice --- the ratings are named in \texttt{VALID\_RAG\_RATINGS} and a division may
  recolour its own bands through the risk policy. These are the platform fallbacks. \\
\texttt{STATE} & Banner and chip tints keyed to the state they announce (danger, warn,
  ok). \\
\texttt{HUE}, \texttt{GRAPH} & Hues that distinguish a state \emph{without} asserting a
  risk rating --- a Product being ``Proposed'' is a lifecycle fact, and drawing it in a RAG
  colour would read as a rating. \\
\texttt{TYPE}, \texttt{SPACE}, \texttt{RADIUS}, \texttt{SHADOW} & The type scale, the
  spacing steps, corner radii and shadows. Both typefaces are tokens too. \\
\bottomrule
\end{longtable}

Two further tables in the same file are about \emph{meaning} rather than value:

\begin{itemize}
  \item \texttt{STATUS\_COLOUR\_TOKENS} --- which \emph{token} each status is drawn in, per
        domain (\texttt{rag}, \texttt{meeting\_status}, \texttt{lifecycle},
        \texttt{product\_status}, \texttt{mrc\_state}, \texttt{validation\_question},
        \texttt{mission\_criticality}, \texttt{review\_thread}). It names a token, never a
        colour. Six maps in \texttt{helpers.js} used to hold 27 entries between them; this
        is those six, unified.
  \item \texttt{STATUS\_COLOUR\_DEFAULTS} --- what an \emph{unrecognised} value is drawn in,
        per domain. A domain absent from it has no default on purpose: an unknown rating
        must render as a dash, not as a grey badge asserting a classification nobody made.
\end{itemize}

Keys are the CSS custom-property name without the leading \texttt{--}, so a token is
spelled the same way in \texttt{config/theme.py}, in the stylesheet, and in an adopter's
override. There is no translation layer to drift.

\subsection{How a token reaches the screen}

\begin{code}
config/theme.py                 the values (THEME, a flat dict of token -> value)
      |
      +--> src/web/_theme.py --> GET /theme.css        :root { --accent: #1976d2; ... }
      |          |                    linked BEFORE app.css, cache-busted by a
      |          |                    SHA-256 of the served text (theme_version)
      |          |
      |          +------------> GET /api/theme         { tokens, status_colours,
      |                              hydrates the JS      status_defaults }
      |                              `Theme` object
      |
      +--> config/reporting.py --> src/governance/latex_style.py --> the PDFs
\end{code}

\texttt{/theme.css} is a \emph{route}, not a checked-in file, for two reasons. A generated
file would have to be regenerated and committed whenever a token changed, and the two would
drift the first time somebody forgot. And a route sidesteps the Jinja template cache: with
\texttt{debug=False} a new \texttt{<script>} or \texttt{<link>} tag in
\texttt{index.html} is not picked up until the server restarts, but the CSS a route serves
is fresh on every request.

Its cache-buster is a hash of the served text rather than the file-modification stamp used
for static assets --- \texttt{asset\_version} watches files under \texttt{src/static}, and
the theme is not one. Hashing gives the same guarantee: the URL changes exactly when the
bytes do, so a rebrand cannot be masked by a browser cache, and an unchanged theme keeps
its cached copy across restarts.

In JavaScript, three accessors on \texttt{Theme} (in \texttt{src/static/js/helpers.js}):

\begin{longtable}{@{}p{0.34\textwidth}p{0.60\textwidth}@{}}
\toprule
\textbf{Call} & \textbf{Use it when} \\
\midrule
\endhead
\texttt{Theme.colour(domain, value)} & Colouring a status --- a rating, a meeting state, a
  lifecycle. The mapping is server-supplied. \\
\texttt{Theme.ref(token)} & You need a reference to a token: emits
  \texttt{var(--token)}. \\
\texttt{Theme.hex(token)} & Only where the value is \emph{manipulated} rather than used ---
  the heatmap appends two hex digits of alpha, and \texttt{var(--red)22} means nothing. \\
\bottomrule
\end{longtable}

\subsection{How a token reaches the PDFs}

\texttt{config/reporting.py} derives the document colours from the theme:

\begin{code}
DOCUMENT_LINK_COLOUR           = BRAND["accent"]   # hyperlinks, headings
DOCUMENT_CHROME_COLOUR         = TEXT["muted"]     # running header/footer furniture
DOCUMENT_CLASSIFICATION_COLOUR = RAG["red"]        # the CONFIDENTIAL banner
DOCUMENT_UNRATED_COLOUR        = RAG["grey"]       # a band with no colour of its own
\end{code}

\texttt{src/governance/latex\_style.py} converts them for \texttt{xcolor}'s HTML model and
applies them; \texttt{latex.py} builds structure only. So a report and the screen it came
from cannot disagree. The same file also owns the page setup taken from the house template
--- geometry, line spacing, rule width, the title fonts --- and
\texttt{DOCUMENT\_BODY\_FONT}, which is empty by default (Computer Modern, the serif face
every generated document has printed in to date) and can be set to \texttt{"sans"} to load
Helvetica and match the console. That is a visible change to every document the platform
produces, so it is a decision an adopter makes rather than a side effect.

The institution's \emph{words} on those documents --- the classification banner, the
``prepared by'' line, the legal notice, the organisation name --- are settings, recorded in
\textbf{Admin \(\rightarrow\) Misc}, and re-label every document at once. The values in
\texttt{config/reporting.py} are what a fresh install prints until it does.

\subsection{Rebranding: the procedure}

\begin{enumerate}
  \item \textbf{Edit \texttt{config/theme.py}.} For most adopters this is the
        \texttt{BRAND} group and nothing else --- \texttt{accent} plus its four variants
        and the three header values. Keep the token \emph{names}; change the values.
  \item \textbf{If the type or spacing scale changes,} edit \texttt{TYPE}, \texttt{SPACE},
        \texttt{RADIUS}, \texttt{SHADOW} in the same file. Values are strings, deliberately
        --- a token is whatever CSS accepts in that position, so \texttt{"13px"},
        \texttt{"0 1px 3px rgba(0,0,0,.12)"} and a whole font stack are all legal.
  \item \textbf{If a status should be drawn in a different token,} edit
        \texttt{STATUS\_COLOUR\_TOKENS} --- not a map in a \texttt{.js} file.
  \item \textbf{Upload the organisation's logo} through the console (Admin); the logo is
        user-uploaded, not a file in the repository, and it flows into the generated
        title pages.
  \item \textbf{Set the document wording} in Admin \(\rightarrow\) Misc: organisation name,
        classification banner, author line, legal notice.
  \item \textbf{Run the styling audit.} \texttt{python -m audit} and read the final line.
        Three gated checks sit at zero: a colour literal in an asset, a colour literal in
        Python, and a \texttt{var(--token)} naming a token the theme does not define.
  \item \textbf{Restart the server and look at it.} The theme route is not cached, but
        \texttt{index.html} is under \texttt{debug=False}, so restart --- or
        \texttt{kill -HUP} the gunicorn master to recycle its workers --- if you touched a
        template. Verify on the scratch database at port 5016
        (\texttt{scripts/restore\_scratch.sh}, then \texttt{scripts/dev\_app.sh}), never
        against the live estate.
  \item \textbf{Regenerate one PDF} and check the brand followed. The Audit Reports tab
        generates them in-app; \texttt{python -m audit --pdf} does it from the command line.
\end{enumerate}

\begin{note}[Three traps, all of which have caught somebody]
\textbf{An undefined token fails silently.} A \texttt{var(--token)} the theme does not
define is not an error --- the browser drops the declaration and the element renders with
whatever it inherits. A rename therefore reaches a screen as a colour quietly falling back
on whichever control nobody happened to open. This is why the audit checks token references
as well as literals, and why you rename by adding the new token and migrating to it rather
than renaming in place.

\textbf{SVG presentation attributes do not accept \texttt{var()}.}
\texttt{fill="var(--accent)"} renders black. Anything drawing into a canvas uses a CSS class
or a \texttt{style} attribute.

\textbf{Many UI surfaces are hand-curated, not spec-driven.} Adding a \emph{field} to
config is not enough --- it has to be wired into the relevant detail lists, the Admin tab
bar and the register form. Colours are the exception, not the pattern: they genuinely do
flow from one file.
\end{note}

\subsection{What is not built yet}

\texttt{docs/design/theme\_config.md} \S4.3 proposes two ways for an adopter to override the
theme without editing code. Neither is built at the commit on the title page:

\begin{itemize}
  \item \textbf{A theme file} --- \texttt{THEME} shipped as defaults, overridden by JSON at
        \cfg{MODELRISK\_THEME}\allowbreak\texttt{=themes/<adopter>.json}. Designed; the
        environment variable is not read anywhere in the code today.
  \item \textbf{Admin \(\rightarrow\) Branding} --- stored like the MRC settings. Deferred
        deliberately.
\end{itemize}

So a rebrand today is a one-file source change, committed and deployed like any other ---
which for a vendor cloning the platform into their own suite is usually acceptable, and is
in any case a change the audit will hold honest. If you need the JSON override, it is a
small piece of work: read the file in \texttt{config/theme.py} after the groups are
assembled, merge over \texttt{THEME}, and leave everything downstream untouched.

\clearpage
\section{The coding rules R1--R10}
\label{sec:rules}

\subsection{Why an integrator should care}

The rules are not house style. Each of the four jobs in this guide is cheap or expensive
depending on whether a particular rule held:

\begin{itemize}
  \item Rebranding is a one-file change \emph{because} \rulename{R10} holds.
  \item Pointing the platform at a different database is configuration \emph{because}
        \rulename{R1} holds; porting it to a different \emph{engine} is bounded \emph{because}
        \rulename{R6} holds.
  \item Containerising is straightforward \emph{because} \rulename{R1} puts every parameter
        behind an environment variable.
  \item Putting your own SSO in front of it is a header-name change \emph{because}
        \rulename{R1} holds for identity too.
\end{itemize}

The full text is \texttt{docs/rules/coding\_rules.md}; each rule there carries the rule,
why it exists, how to comply, and how it is verified. This section is the operator's
summary.

\subsection{The rules}

\begin{longtable}{@{}p{0.05\textwidth}p{0.42\textwidth}p{0.45\textwidth}@{}}
\toprule
& \textbf{Rule} & \textbf{Enforcement} \\
\midrule
\endhead
R1 & No hardcoded parameters outside \texttt{config/}. Filenames, ports, header names,
     field lists, value sets, thresholds, default modes. Derive, don't duplicate. &
     \textbf{Gated} by the hardcoding scan. \\
R2 & No non-test source file over \textbf{300} lines; test files over \textbf{400}.
     Split into a package of focused modules. &
     \textbf{Gated} for \texttt{.py} repo-wide (\texttt{config}, \texttt{src},
     \texttt{scripts}, \texttt{app}, \texttt{docs} plus root files) and for tests. JS/CSS
     are reported, not gated. \\
R3 & \(\geq\)\textbf{99\%} line+branch coverage, checked every stage. &
     \textbf{Gated} by \texttt{fail\_under = 99} in \texttt{pyproject.toml}; the audit
     prints the current figure and the weakest packages. \\
R4 & No functions in \texttt{\_\_init\_\_.py} --- imports and \texttt{\_\_all\_\_} only. &
     \textbf{Gated}. The audit lists every package initialiser with its length and its
     re-export verdict. \\
R5 & Every first-party file carries the canonical licence header or a one-line notice;
     files that cannot hold a comment get a sibling \texttt{.license}. &
     \textbf{Deliberately not automated.} Whose name belongs at the top of a file is a
     commercial question, and a failing build forces it at the wrong moment. The
     maintainer runs a local-only tool. \\
R6 & Never touch the data store directly. All persistence through
     \texttt{src/database/}; no SQL, cursor, connection or driver import anywhere else.
     No escape hatch. &
     \textbf{Reported}, in three scans: data access outside the package, path definitions
     outside \texttt{config}, and \texttt{.json} file I/O under the source tree. Treated as
     zero-tolerance in review. \\
R7 & Front-end assets under \texttt{src/static} and \texttt{src/templates}; never JS or
     CSS embedded in a \texttt{.py} file. &
     \textbf{Gated} for embedded JS/CSS; the directory layout is reviewed, and a misplaced
     asset shows up as a 404 in the e2e smoke suite. \\
R8 & No duplication --- extract a named utility. Zero tolerance across two models. &
     \textbf{Reported}: jscpd's duplication percentage (threshold 5\%, clones of
     \(\geq\)8 lines / 50 tokens). Cross-model duplication is a hard fault in review. \\
R9 & Test and verify against a restored backup, never the live database. &
     \textbf{Partially enforced.} \texttt{assert\_disposable} guards the destructive test
     fixtures; \texttt{scripts/dev\_app.sh} refuses to start against a protected database.
     An ordinary write to a protected database is not stopped --- that gap is what the
     known-actor rule (\S\ref{sec:permissions}) narrows. \\
R10 & Every visual parameter comes from \texttt{config/theme}. &
     \textbf{Gated} on three checks, all at zero: a colour literal in an asset, a colour
     literal in Python, and a \texttt{var(--token)} naming an undefined token. \\
\bottomrule
\end{longtable}

\subsection{Running the audit}

\begin{code}
python -m audit           # console report; exits non-zero on a gated finding
python -m audit --pdf     # also renders one PDF per rule plus a consolidated overview
\end{code}

The last line is the verdict:

\begin{code}
==================================================
PASS: 0 gated finding(s)
\end{code}

\begin{note}[Read the line, not the exit code]
A background run that reports ``completed, exit code 0'' with empty output is \textbf{not}
a pass --- that exit code belongs to the shell that launched it. Always read the final
\texttt{PASS:} / \texttt{FINDINGS:} line.
\end{note}

Sections marked \texttt{OK} have no findings and \texttt{!!} have some; the reported
(non-gating) sections are printed after the gated ones and are not counted in the total.
Coverage is read from \texttt{coverage.xml} if one is present, so for a figure rather than a
reminder, run \texttt{coverage run -m pytest \&\& coverage xml} first. Duplication needs
Node and \texttt{jscpd}; without them that line says so and the audit continues.

\subsection{The reports}

\texttt{--pdf} --- and the console's \textbf{Audit Reports} tab, which calls the same
function so both emit exactly the same set --- writes one PDF per rule plus a consolidated
overview:

\begin{code}
audit_r1_hardcoded_parameters.pdf     audit_r7_embedded_js.pdf
audit_r2_large_files.pdf              audit_r8_duplicate_code.pdf
audit_r3_test_coverage.pdf            audit_r10_styling.pdf
audit_r4_init_files.pdf               json_files_report.pdf
audit_r6_data_access.pdf              coding_rule_audit.pdf   (the overview)
audit_r6_path_definitions.pdf
\end{code}

They are written to \texttt{AUDIT\_OUTPUT\_DIR}, which is
\texttt{\$UPLOAD\_DIR/audit/<database name>} --- under the configured upload directory, so
generated state follows \cfg{UPLOAD\_DIR} wherever an operator points it rather than being
written into the checkout, and scoped by database so reports from a demo and a live
instance never masquerade as one another. Rendering needs tectonic; without it the audit
says so and returns non-zero for that reason alone.

Alongside them the audit run drops its test evidence --- \texttt{junit.xml},
\texttt{coverage.xml}, the browsable HTML coverage report, and
\texttt{init\_audit\_results.json} --- so one directory holds the whole quality story for a
given database.

\subsection{Where the thresholds live}

Every number in this section is a value in \texttt{config/audit.py}, not a literal in the
audit code (\rulename{R1} applies to the audit too):
\texttt{MAX\_SOURCE\_FILE\_LINES = 300}, \texttt{MAX\_TEST\_FILE\_LINES = 400},
\texttt{COVERAGE\_MINIMUM = 99}, \texttt{DUPLICATION\_MAX\_PERCENT = 5.0},
\texttt{JSCPD\_MIN\_LINES = 8}, \texttt{JSCPD\_MIN\_TOKENS = 50}. The scan scope is there
too: \texttt{SOURCE\_DIRS}, \texttt{REPO\_SCAN\_DIRS}, \texttt{STYLING\_ASSET\_DIRS},
\texttt{STYLING\_GATED\_SUFFIXES}, \texttt{EXCLUDE\_DIRS}.

An adopting institution that wants a different bar changes those values. Be aware that
loosening one is visible in the generated report --- the PDFs print the threshold they were
run against, which is the point.

\subsection{The baseline}

A rule set is only useful if ``clean'' means something. At the commit on the title page the
gate sits at \textbf{2 findings}, both R2, both known:

\begin{code}
384  src/governance/client/pack.py
326  src/web/api_mrc/lifecycle.py
\end{code}

Anything beyond those two is new work's doing. Diff against \texttt{main} before assuming
otherwise --- the list moves as files are split.

\clearpage
\section{Switching to a different SQL server}
\label{sec:sqlserver}

\subsection{Two different questions}

``Switching the SQL server'' means one of two things, and they are days apart in effort.

\paragraph{A different PostgreSQL server --- configuration only.} A managed instance (RDS,
Cloud SQL, Azure Database for PostgreSQL), a different host, port, database name or role
set is \textbf{no code change at all}. Seven environment variables, read in
\texttt{config/database.py}:

\begin{code}
DB_HOST=db.internal.example.com
DB_PORT=5432
DB_NAME=modelrisk
DB_OWNER_USER=modelrisk_owner        # owns the schema, runs migrations
DB_OWNER_PASSWORD=...
DB_APP_USER=modelrisk_app            # the runtime role; append/read only on events
DB_APP_PASSWORD=...
# DB_MAINTENANCE_NAME=postgres       # only used while creating a database
\end{code}

Start the app; \texttt{bootstrap()} applies the migrations, seeds the admin reference lists
on a fresh database, and rebuilds the read models. Two constraints on the names you choose:
\cfg{DB\_APP\_USER} must be a plain SQL identifier (the migrations grant to it by name and
the runner refuses anything else), and \cfg{DB\_NAME} must not collide with a name in
\texttt{PROTECTED\_DB\_NAMES} unless it really is the estate --- that list is what the test
fixtures and the verification app refuse to touch.

\paragraph{A different engine --- a port.} Microsoft SQL Server, Oracle, MySQL. That is
real work, and the rest of this section is about scoping it honestly. \textbf{Nothing in
the repository implements a second engine today.} What exists is
\texttt{docs/mssql-immutability.md}, which covers the one part of a SQL Server port that is
easy to get wrong and impossible to notice.

\subsection{What is engine-specific, and where}

The good news is that rule \rulename{R6} makes the answer enumerable. Every module outside
\texttt{src/database/} reaches the store through named functions, so nothing above that
line knows what engine it is talking to. The engine-specific surface is:

\begin{longtable}{@{}p{0.34\textwidth}p{0.60\textwidth}@{}}
\toprule
\textbf{Where} & \textbf{What is PostgreSQL-specific} \\
\midrule
\endhead
\texttt{src/database/\_connection.py} & The driver. \texttt{psycopg} imports, the
  \texttt{connect()} calls, the autocommit admin connection, and
  \texttt{CREATE DATABASE} via \texttt{psycopg.sql.Identifier}. This is the only module
  that opens a connection. \\
\texttt{src/database/eventstore.py} & The SQL itself. \texttt{\%s} parameter markers,
  \texttt{INSERT ... RETURNING}, and \texttt{psycopg.types.json.Jsonb} for the payload
  column. \\
\texttt{src/database/\_migrate.py} & \texttt{pg\_advisory\_xact\_lock} for the
  one-migrator-at-a-time guarantee, and \texttt{TIMESTAMPTZ}/\texttt{now()} in the
  \texttt{schema\_migrations} DDL. \\
\texttt{src/database/}\newline\texttt{migrations/*.sql} & Eleven files of PostgreSQL DDL:
  \texttt{BIGINT GENERATED ALWAYS AS IDENTITY}, \texttt{UUID} with
  \texttt{gen\_random\_uuid()}, \texttt{JSONB}, \texttt{TIMESTAMPTZ}, the
  \texttt{DO \$\$ ... \$\$} block that creates the role, and \texttt{GRANT}/\texttt{REVOKE}
  semantics. \\
\texttt{src/database/projections.py} & Read-model upserts (\texttt{ON CONFLICT}) and any
  JSON operators used to query the payload. \\
\texttt{scripts/pg.sh}, \texttt{scripts/backup\_db.sh}, \texttt{restore\_scratch.sh} &
  \texttt{pg\_ctl}, \texttt{pg\_dump}, \texttt{psql}. Operational, not application code, but
  a deployment without a working backup is not a deployment. \\
\bottomrule
\end{longtable}

Everything else --- roughly the whole application --- is untouched. That is the dividend
R6 was written to pay.

\subsection{A port, step by step}

\begin{enumerate}
  \item \textbf{Do not add T-SQL to \texttt{src/database/migrations/}.} The runner globs
        \texttt{*.sql} and applies \emph{every} file in lexical order as PostgreSQL; a
        T-SQL file dropped in there is executed against Postgres and fails the migration for
        every existing deployment. Give the new engine its own migrations directory and
        select it in \texttt{config/paths.py}.
  \item \textbf{Introduce a backend seam inside \texttt{src/database/}.} Keep the public
        names --- \texttt{append\_event}, \texttt{read\_stream}, \texttt{read\_view},
        \texttt{replay}, \texttt{run\_migrations} --- and put the dialect behind them. A
        \cfg{DB\_BACKEND} variable in \texttt{config/database.py} selects it (\rulename{R1}:
        the choice is a parameter, not an import somebody edits).
  \item \textbf{Translate the schema.} \texttt{IDENTITY} for \texttt{global\_seq},
        \texttt{UNIQUEIDENTIFIER}/\texttt{NEWID()} or an application-generated UUID for
        \texttt{event\_id}, \texttt{NVARCHAR(MAX)} with an \texttt{ISJSON} check for the
        payload, \texttt{DATETIMEOFFSET} for the timestamps. Keep the unique constraint on
        \texttt{(aggregate\_type, aggregate\_id, stream\_seq)}: it is the concurrency
        backstop, not an index.
  \item \textbf{Replace the advisory lock.} On SQL Server, \texttt{sp\_getapplock} with
        \texttt{@LockOwner = 'Transaction'} is the direct equivalent. Do not skip it because
        ``we only run one process'' --- gunicorn forks, and this is exactly the failure that
        looks like a worker crashing on boot for no reason.
  \item \textbf{Re-establish immutability at the grant level.} This is the step that is
        easy to lose, and it has its own document. See below.
  \item \textbf{Fix the backups.} \texttt{scripts/backup\_db.sh} locates \texttt{pg\_dump}
        and \textbf{exits 0 when it cannot find one}, so on another engine the hourly
        schedule reports success while doing nothing. Replace it, and verify a restore, not
        just a dump.
  \item \textbf{Run the suite against the new engine.} \texttt{pytest} derives its own
        disposable database and rebuilds every projection, so a passing suite is a genuine
        statement about the port. Then \texttt{python -m audit} --- a port that leaks a
        driver import outside \texttt{src/database/} shows up in the R6 section.
\end{enumerate}

\subsection{Immutability does not cross the border}

The platform's central claim is that the event log is the system of record and cannot be
rewritten. On PostgreSQL that is enforced by the database:

\begin{code}
GRANT SELECT, INSERT ON events TO {{app_user}};
REVOKE UPDATE, DELETE, TRUNCATE ON events FROM {{app_user}};
\end{code}

A driver swap moves the queries. It does not move the grants --- those are DDL that the
PostgreSQL migration applies. Until the equivalent is in place on the new engine, the
runtime account can rewrite the log: the hash chain will still \emph{detect} it, but nothing
\emph{prevents} it, and ``enforced at the database-grant level'' is a claim you may be
making to an assessor.

On SQL Server there is one difference that turns a correct-looking script into a decorative
one. \textbf{\texttt{REVOKE} only removes a permission granted directly}; it does not
override a permission inherited through role membership. If the application user is a member
of \texttt{db\_datawriter} --- a very common way to set up an app account ---
\texttt{REVOKE UPDATE ON events} changes nothing at all and the log stays writable.
\texttt{DENY} is the one that wins. Two smaller differences: there is no \texttt{TRUNCATE}
permission (it is authorised by \texttt{ALTER}, so you \texttt{DENY ALTER}), and
\texttt{sysadmin} bypasses permission checks entirely --- as will \texttt{db\_owner} ---
so check what the app connects as before trusting any of it.

\begin{code}
GRANT SELECT, INSERT ON dbo.events TO [mkm_app];
DENY UPDATE, DELETE ON dbo.events TO [mkm_app];
DENY ALTER ON dbo.events TO [mkm_app];      -- TRUNCATE is authorised by ALTER
\end{code}

The projections stay mutable --- they are derived and rebuildable, exactly as on
PostgreSQL. \texttt{docs/mssql-immutability.md} carries the full script, the two privilege
checks that prove the account is not privileged enough to ignore the \texttt{DENY}, and the
functional test that an assessor should actually accept:

\begin{code}
EXECUTE AS LOGIN = N'mkm_app';
BEGIN TRY
    UPDATE dbo.events SET actor = actor WHERE 1 = 0;
    PRINT '*** FAILED: the event log is writable ***';
END TRY
BEGIN CATCH
    PRINT 'Blocked as expected: ' + ERROR_MESSAGE();
END CATCH
REVERT;
\end{code}

\texttt{WHERE 1 = 0} matches no rows --- SQL Server checks the permission on the object
regardless, so the test cannot alter data even if the \texttt{DENY} were missing. Expect
error 229, \emph{the UPDATE permission was denied on the object 'events'}.

\begin{note}[Two things immutability grants do not cover]
The \textbf{owner / migration account can still alter the log}, by design --- it has to be
able to create the schema. Immutability is a property of the \emph{runtime} account, which
is why the two roles are separate configuration values and why the owner credentials should
not be in the application's environment beyond the moment migrations run.

The \textbf{hash chain is the detection half} and is application-level, so it crosses to any
engine unchanged. The grants are the prevention half. An audit story wants both.
\end{note}

\clearpage
\section{Running it inside a container}
\label{sec:container}

\subsection{Status: not shipped}

There is \textbf{no \texttt{Dockerfile} or Compose file in the repository} at the commit on
the title page. The platform is developed and deployed natively --- the README says so
explicitly (``PostgreSQL 16, run natively, no Docker''). The listings below are a worked
starting point that follows the application's actual startup contract; they are not
something the test suite exercises, so treat them as a design you own rather than as a
supported artefact.

The application itself is container-friendly for one reason: \rulename{R1}. Every parameter is
already an environment variable with a shipped default, so there is nothing to bake into an
image and no file to template.

\subsection{What has to be in the image}

\begin{longtable}{@{}p{0.26\textwidth}p{0.14\textwidth}p{0.54\textwidth}@{}}
\toprule
\textbf{Component} & \textbf{Required?} & \textbf{Why} \\
\midrule
\endhead
Python 3.13+ & Yes & \texttt{requires-python = ">=3.13"}. \\
\texttt{psycopg[binary]}, \texttt{flask}, \texttt{gunicorn}, \texttt{jsonschema} & Yes &
  The runtime dependencies. \texttt{gunicorn} needs \texttt{fcntl}, so it is POSIX-only ---
  on Windows the development server is the only option. \\
\textbf{tectonic} & Effectively & The \LaTeX{} engine. Without it every PDF endpoint ---
  meeting pack, Model Risk Report, audit reports --- returns \texttt{503}. For a governance
  platform whose output is committee papers, that is not really optional. \\
pandoc & Optional & Only the User Guide PDF in the Audit Reports tab. \\
\texttt{postgresql-client} & Conditional & \texttt{scripts/backup\_db.sh} needs
  \texttt{pg\_dump}. Omit it \emph{and} set \cfg{BACKUP\_ENABLED=0} --- see below. \\
Node + \texttt{jscpd} & No & Only the duplicate-code section of \texttt{python -m audit};
  keep it out of a runtime image. \\
\bottomrule
\end{longtable}

\subsection{A worked Dockerfile}

\begin{code}
FROM python:3.13-slim AS base

# tectonic renders every PDF the platform produces; pandoc is only the user guide.
# postgresql-client provides pg_dump for the in-app backup schedule -- drop both it
# and BACKUP_ENABLED if the database is backed up by the platform hosting it.
RUN apt-get update && apt-get install -y --no-install-recommends \
        tectonic pandoc postgresql-client ca-certificates \
    && rm -rf /var/lib/apt/lists/*

WORKDIR /app

COPY requirements.txt ./
RUN pip install --no-cache-dir -r requirements.txt

COPY . .

# Warm tectonic's package cache at BUILD time, so the first PDF request in
# production does not have to reach the internet. Without this a locked-down
# network turns the first report into a timeout rather than a 503.
ENV TECTONIC_CACHE_DIR=/opt/tectonic-cache
RUN mkdir -p "$TECTONIC_CACHE_DIR" \
    && printf '%s\n' \
        '\documentclass{article}' \
        '\usepackage{xcolor,longtable,booktabs,geometry,fancyhdr,hyperref}' \
        '\begin{document}warm\end{document}' > /tmp/warm.tex \
    && tectonic -o /tmp /tmp/warm.tex && rm -f /tmp/warm.*

# Run as a non-root user that owns the state directory.
RUN useradd --system --uid 10001 --home /app modelrisk \
    && mkdir -p /data/uploads /opt/tectonic-cache \
    && chown -R modelrisk:modelrisk /data /opt/tectonic-cache
USER modelrisk

# The container's network namespace is not the host's: bind the wildcard address,
# not the 127.0.0.1 that config/web.py defaults to.
ENV WEB_HOST=0.0.0.0 \
    WEB_PORT=5015 \
    UPLOAD_DIR=/data/uploads \
    PYTHONUNBUFFERED=1

EXPOSE 5015
VOLUME ["/data/uploads"]

# /api/models proves two things at once: the process serves, and it can reach
# the database. There is no dedicated /health route. The start period covers
# migrations and the read-model replay, which run before the first request.
HEALTHCHECK --interval=15s --timeout=10s --start-period=60s --retries=5 \
  CMD python -c "import urllib.request as u; \
      u.urlopen('http://127.0.0.1:5015/api/models')" || exit 1

# gunicorn.conf.py runs bootstrap() in the master before the fork: migrations,
# seeded reference lists, read-model rebuild. Workers inherit the done-marker.
CMD ["gunicorn", "--config", "gunicorn.conf.py", "wsgi:application"]
\end{code}

\subsection{A worked Compose file}

\begin{code}
services:
  db:
    image: postgres:16
    environment:
      POSTGRES_USER: mkm_owner            # must equal DB_OWNER_USER below
      POSTGRES_PASSWORD: ${DB_OWNER_PASSWORD:?set it}
      POSTGRES_DB: mkm_modelrisk
    # The app opens workers x threads connections; the stock max_connections is 100.
    command: ["postgres", "-c", "max_connections=200"]
    volumes:
      - pgdata:/var/lib/postgresql/data
    healthcheck:
      test: ["CMD-SHELL", "pg_isready -U mkm_owner -d mkm_modelrisk"]
      interval: 5s
      timeout: 3s
      retries: 20

  app:
    build: .
    depends_on:
      db:
        condition: service_healthy       # migrations run at boot; the DB must be up
    environment:
      DB_HOST: db
      DB_PORT: 5432
      DB_NAME: mkm_modelrisk
      DB_OWNER_USER: mkm_owner
      DB_OWNER_PASSWORD: ${DB_OWNER_PASSWORD:?set it}
      DB_APP_USER: mkm_app               # created by migration 0002 with these grants
      DB_APP_PASSWORD: ${DB_APP_PASSWORD:?set it}
      IDENTITY_HEADER: X-Auth-Request-User
      ENFORCE_KNOWN_ACTORS: always       # see section 7
      IDENTITY_DEV_FALLBACK: unauthenticated   # deliberately not a registered person
      BACKUP_ENABLED: 0                  # back the volume/database up outside the app
      WEB_WORKERS: 4
      WEB_THREADS: 4
      WEB_TIMEOUT: 300                   # tectonic renders take far longer than 30s
    volumes:
      - uploads:/data/uploads
    ports:
      - "127.0.0.1:5015:5015"            # never publish this straight to a network
    # The health check is declared in the Dockerfile (HEALTHCHECK) and inherited.

volumes:
  pgdata:
  uploads:
\end{code}

\subsection{The things that only bite once it is in a container}

\paragraph{Bind the wildcard address.} \texttt{config/web.py} defaults
\cfg{WEB\_HOST} to \texttt{127.0.0.1}, which inside a container means ``reachable by nothing''.
Set \cfg{WEB\_HOST=0.0.0.0} (or \cfg{WEB\_BIND} directly) and put the access control at the
proxy, not at the bind address.

\paragraph{Boot is not instant, and it touches the database.} The gunicorn master applies
migrations, seeds reference lists on a fresh database, and \emph{replays every projection}
before it forks. On a large estate that is the slowest thing the process ever does. Give the
health check a \texttt{start\_period}, make the app wait for the database rather than
restart-looping, and if replay-per-boot becomes unacceptable set
\cfg{REBUILD\_VIEWS\_ON\_STARTUP=0} --- knowing that you have then taken responsibility for
rebuilding a view after a restore.

\paragraph{One master, many workers --- do not scale the container blindly.} Two things are
deliberately once-per-boot in the master: the bootstrap and the backup schedule. Running
three replicas gives you three backup schedules and three processes racing the same
migrations. The advisory lock makes the race safe, not free. If you need horizontal scale,
run one instance with the schedule enabled and the rest with \cfg{BACKUP\_ENABLED=0}, or
lift bootstrap into an init container / \texttt{Job} and start the app with migrations
already applied.

\paragraph{Size the pool against \texttt{max\_connections}.} Peak connections are
\texttt{WEB\_WORKERS \(\times\) WEB\_THREADS} --- 16 by default --- opened and closed per
call with no shared pool. The defaults are deliberately modest rather than the usual
\((2 \times \mathrm{cores}) + 1\), because that formula on a many-core host exhausts a stock
\texttt{max\_connections} of 100.

\paragraph{PDFs need time and a warm cache.} \cfg{WEB\_TIMEOUT} is 300 seconds for exactly
one reason: rendering shells out to tectonic, and under gunicorn's stock 30-second timeout
the worker is killed mid-render and the request fails as a bare \texttt{502} with nothing in
the log to explain it. Tectonic also fetches \LaTeX{} packages on first use --- warm the
cache at build time (above) or the first report in a network-restricted environment hangs
rather than failing cleanly.

\paragraph{Uploads are state; the container is not.} \cfg{UPLOAD\_DIR} holds every uploaded
document and every generated report (audit PDFs land in
\texttt{\$UPLOAD\_DIR/audit/<database name>}). It must be a volume, and it must be in your
backup set --- the event log records that a document exists and its metadata, not its bytes.

\paragraph{Do not bake \texttt{.env} into the image.} \texttt{config/env.py} loads
\texttt{.env} at import if it is present. A real environment variable always wins, so an
\texttt{.env} that got copied in during \texttt{COPY . .} will not override your
orchestrator's settings --- but it \emph{will} fill in anything you forgot to set, silently,
including a \cfg{DB\_NAME} pointing at somebody else's estate. Add it to
\texttt{.dockerignore} and keep secrets in the platform's secret store.

\paragraph{Backups inside a container are a decision, not a default.}
\texttt{scripts/backup\_db.sh} sources \texttt{.env} for its own credentials and shells out
to \texttt{pg\_dump}; if it cannot find one it exits 0, so a misconfigured image reports a
healthy hourly schedule while writing nothing. In a container, prefer the database
platform's own snapshots and set \cfg{BACKUP\_ENABLED=0} --- or keep the schedule, install
\texttt{postgresql-client}, mount a backup volume, and \emph{verify a restore}.

\paragraph{Verify against a scratch copy, not the estate.} Rule \rulename{R9} does not relax
because the app is in a container. A verification instance points \cfg{DB\_NAME} at a
scratch database restored from a backup; the guard in \texttt{scripts/dev\_app.sh} that
refuses to start against a protected name is a shell script, so a container that invokes
gunicorn directly does not inherit it. Either keep that entry point or re-implement the
check in yours.

\clearpage
\section{Permissions and SSO}
\label{sec:permissions}

\subsection{The model in one sentence}

\textbf{The platform does not authenticate anybody.} Identity is established at the
perimeter --- by your SSO, a reverse proxy, an identity-aware gateway --- and forwarded to
the application in a request header, which the application \emph{trusts verbatim} and
records as the author of every event it writes.

That is a legitimate architecture, and it is the same one oauth2-proxy, Google IAP, Azure
Application Gateway and most Kubernetes ingress-auth setups assume. It has exactly one
requirement, and \S\ref{sec:trust} is about taking it seriously: \textbf{if anything can
reach the application port without passing through the proxy, it can author governance
history under any name.}

\subsection{What arrives, and what happens to it}

\begin{code}
config/identity.py
  IDENTITY_HEADER       = "X-Auth-Request-User"                  (the stable identifier)
  IDENTITY_NAME_HEADER  = "X-Auth-Request-Preferred-Username"    (display name, optional)
  IDENTITY_DEV_FALLBACK = "<an address>"   used when NO header is present
  SYSTEM_ACTOR          = "system"         the platform's own actor, for first-run seeding
\end{code}

\texttt{resolve\_principal()} (\texttt{src/identity/resolver.py}) does a case-insensitive
header lookup and returns a \texttt{Principal} carrying \texttt{actor},
\texttt{display\_name} and a \texttt{source} --- \texttt{sso\_header}, \texttt{fallback} or
\texttt{ingest}. It is a deliberate seam: the rest of the application needs only a
\texttt{Principal}, so a different identity source can be substituted here without touching
a single command handler.

Machine producers get a synthetic principal, \texttt{ingest:<source\_repo>}, so an ingested
payload can never be mistaken in the audit trail for a person's work.

\subsection{Who is allowed to author a change}

Phase P1 of \texttt{docs/permissions\_design.md} is \textbf{built and enforced}: a change
must be authored by somebody the institution has registered under
\textbf{Admin \(\rightarrow\) People}. The check lives in \texttt{authorize()} in
\texttt{src/command/dispatch.py} --- so every mutating command passes it, and a new endpoint
cannot forget to ask --- and in \texttt{request\_actor()} for the upload paths that call
their command directly. A refusal is an HTTP \texttt{403} naming the actor and the remedy.

Matching is on \textbf{name or email}, case- and whitespace-insensitive, because the console
sends a person's name and an SSO deployment sends their address; either identifies them.

Policy is \texttt{config/permissions.py}:

\begin{longtable}{@{}p{0.34\textwidth}p{0.60\textwidth}@{}}
\toprule
\textbf{Setting} & \textbf{Meaning} \\
\midrule
\endhead
\cfg{ENFORCE\_KNOWN\_ACTORS} \newline \texttt{protected} (default) &
  The rule applies only against a protected database --- the real estate. Test and scratch
  databases stay open so the suite and throwaway environments are unaffected. \\
\hspace{1em}\texttt{always} & Every database. \textbf{What a hardened deployment sets.} \\
\hspace{1em}\texttt{never} & Disabled. A break-glass switch for an adopter mid-migration;
  it leaves the governance record open to any identity, so it should be temporary. \\
\cfg{ENFORCE\_REVIEWER\_}\newline\cfg{INDEPENDENCE} \newline \texttt{always} (default) &
  Segregation of duties: a model's peer reviewer may not be its own owner. Existing
  breaches are reported by the completeness audit either way; enforcement only refuses the
  write that would create or keep one, so no record is ever locked shut. \\
\texttt{MACHINE\_ACTOR\_PREFIX} = \texttt{ingest:} & Machine producers pass the known-actor
  rule without being people. \\
\texttt{ALLOW\_WHEN\_NO\_}\newline\texttt{PEOPLE\_REGISTERED} & While the People list is empty the rule
  stands down, so a fresh install can register its first person. \textbf{The first person
  added closes it.} \\
\texttt{ROLES} & The role vocabulary (owner, validator, committee, chair, administrator,
  reader), with the explanation shown beside each choice in Admin. An adopting institution
  re-sets this to its own framework. \\
\bottomrule
\end{longtable}

\begin{note}[Renaming a role un-assigns it]
\texttt{ROLES} keys are stored on the People records. Changing a key's \emph{name} without
migrating the records that hold it silently un-assigns the role --- exactly as renaming a
risk band would. On a live estate, add the new role and migrate; never rename in place.
\end{note}

\subsection{The trust boundary --- the part to get right}
\label{sec:trust}

The header is taken at face value. On 2026-08-06 the identity hard-coded in the web tests,
\texttt{alice@corp}, authored 24 real events on the live estate from an ad-hoc script, and
the audit trail recorded them faithfully because nothing was wrong with them --- they were
correctly attributed to the identity that was sent. P1 narrows \emph{who} can author (a
registered person) but does not verify that the caller \emph{is} that person: anyone who can
address the app can still send a registered name.

So the deployment has to close it. Five things, all outside the application:

\begin{enumerate}
  \item \textbf{The app listens only where the proxy can reach it.} Loopback, a private
        network, or a Kubernetes \texttt{ClusterIP} --- never a published port on a routable
        interface. The Compose example above binds \texttt{127.0.0.1:5015} for this reason.
  \item \textbf{The proxy strips the identity headers from every inbound request} before
        setting its own. This is the single most important line of configuration in an SSO
        integration, and it is the one most often left out: without it a client simply sends
        \texttt{X-Auth-Request-User} itself.
  \item \textbf{Set \cfg{ENFORCE\_KNOWN\_ACTORS=always}} rather than relying on the database
        name being one the platform happens to recognise as protected.
  \item \textbf{Point \cfg{IDENTITY\_DEV\_FALLBACK} at something that is not a person.}
        When no header arrives the web layer falls back to this actor --- it does not
        refuse. Set it to a value nobody will ever register (\texttt{unauthenticated}), and
        the known-actor rule turns every header-less write into a \texttt{403} instead of
        an event attributed to whoever the default names.
  \item \textbf{Prefer a stable, non-reassignable identifier} in the header. An email
        address that can be reassigned to a new joiner makes historical attribution
        ambiguous years later, which for an append-only governance record is a real problem
        rather than a theoretical one.
\end{enumerate}

\subsection{Wiring it to an identity provider}

The application only needs the header names to match, so any of these work. Set
\cfg{IDENTITY\_HEADER} and \cfg{IDENTITY\_NAME\_HEADER} to whatever your layer emits; the
defaults follow the oauth2-proxy convention.

\paragraph{oauth2-proxy in front of nginx (OIDC --- Entra ID, Okta, Ping, Keycloak).}

\begin{code}
server {
    listen 443 ssl;

    # 1. Never let a client supply its own identity.
    proxy_set_header X-Auth-Request-User                "";
    proxy_set_header X-Auth-Request-Preferred-Username  "";
    proxy_set_header X-Auth-Request-Email               "";

    location /oauth2/ {
        proxy_pass http://oauth2-proxy:4180;
        proxy_set_header X-Auth-Request-Redirect $request_uri;
    }

    location = /oauth2/auth {
        proxy_pass http://oauth2-proxy:4180;
        proxy_pass_request_body off;
        proxy_set_header Content-Length "";
    }

    location / {
        auth_request /oauth2/auth;
        error_page 401 = /oauth2/start?rd=$request_uri;

        # 2. Set them from what the auth subrequest returned.
        auth_request_set $user  $upstream_http_x_auth_request_user;
        auth_request_set $pref  $upstream_http_x_auth_request_preferred_username;
        proxy_set_header X-Auth-Request-User               $user;
        proxy_set_header X-Auth-Request-Preferred-Username $pref;

        proxy_pass http://modelrisk:5015;
        proxy_read_timeout 300s;    # match WEB_TIMEOUT: PDF renders are slow
    }
}
\end{code}

\paragraph{Other layers.} Google IAP forwards
\texttt{X-Goog-Authenticated-User-Email}; Azure App Service Easy Auth forwards
\texttt{X-MS-CLIENT-PRINCIPAL-NAME}; an Istio/Envoy \texttt{ext\_authz} filter emits
whatever you configure. In each case point \cfg{IDENTITY\_HEADER} at it and keep the
strip-then-set discipline. Where the provider emits a signed assertion (a JWT) rather than a
plain header, either have the gateway verify it and emit the plain header, or replace
\texttt{resolve\_principal()} --- that is what the seam is for, and it is the only function
that would change.

\paragraph{Then register the people.} The known-actor rule matches against
\textbf{Admin \(\rightarrow\) People}, so the identities your IdP emits have to appear there,
as each person's name or email. Two practical notes: the very first person can be added
while the list is still empty (the rule stands down until then), and if you provision People
in bulk from a directory export, use the same string form the header will carry.

\subsection{What is designed but not built}

\begin{longtable}{@{}p{0.10\textwidth}p{0.50\textwidth}p{0.34\textwidth}@{}}
\toprule
\textbf{Phase} & \textbf{What it is} & \textbf{Status} \\
\midrule
\endhead
P1 & Know who is real --- a change must be authored by a registered person. &
     \textbf{Built}, 2026-08-07. \\
P2 & Roles as data on the People record; assignment is itself an event. &
     Vocabulary and \texttt{roles\_for()} exist; nothing consumes them but the Admin
     screen that sets them. \\
P3 & \texttt{authorize()} implemented in \emph{report-only} mode --- log refusals for a
     week of real use. & Designed. This is the phase that stops a permissions rollout
     blocking the day job. \\
P4 & Enforce: refusals become \texttt{403} naming the missing role. & Designed. \\
P5 & Segregation of duties --- an owner cannot approve their own model; a validator cannot
     validate one they own; an administrator cannot ratify the policy they drafted. &
     Reviewer independence is enforced; the rest is designed. \\
P6 & Machine actors scoped to append-only. & Designed. \\
\bottomrule
\end{longtable}

The design's organising principle is worth stating, because it constrains how you should
extend it: \textbf{permissions come from the domain, not a parallel user store}. Who is a
real person is the People list; who owns a model is the model's owner field; who sits on the
committee that governs it is the MRC's membership. There is no user table and no second
role store to keep in step --- which also means permissions are themselves auditable and
ratifiable, exactly like everything else in the record.

When you do implement P3--P5, the permission map belongs in \texttt{config/permissions.py}
as data (\rulename{R1}) --- \texttt{COMMAND\_ROLES = \{"update\_field": ("owner",
"administrator"), "approve\_signoff": ("committee",), ...\}} --- and the enforcement belongs
in \texttt{authorize()}, which every mutating command already passes through. The UI should
stop offering what the caller may not do, but that is cosmetic: the server is the control.

\clearpage
\appendix

\section{Feeding an ERM}
\label{sec:erm}

\subsection{What this surface is for}

Model risk is rarely the end of the story. Where an institution runs enterprise risk
management, model risk feeds two places in it --- the \textbf{operational risk} view and
\textbf{strategic direction} --- and those are read by different people asking different
questions. The API is shaped around that rather than around how this platform happens to
store things.

\begin{code}
GET /api/erm/v1/model-risk             the panel: verdict, operational, direction, estate
GET /api/erm/v1/model-risk/attention   the queue: every outstanding obligation
GET /api/erm/v1/model-risk/estate      the position: counts by kind/tier/rating/division
\end{code}

All three accept \texttt{?as\_of=YYYY-MM-DD}, which moves the dated rules (review currency,
remediation dates). It does \emph{not} replay history: the registers hold current state.

\textbf{The ERM calls you. It never writes.} There is no write endpoint on this surface at
any version, and there will not be one: if the ERM could change the governance record, the
record would no longer be the thing the committee signed.

\subsection{The blocks, and who reads them}

\begin{longtable}{@{}p{0.20\textwidth}p{0.28\textwidth}p{0.46\textwidth}@{}}
\toprule
\textbf{Block} & \textbf{Consumed by} & \textbf{Carries} \\
\midrule
\endhead
\texttt{verdict} & the C-suite tile & Whether anything needs a response, as a sentence and
  a count per state. \\
\texttt{operational} & operational risk & Failures that actually happened in the rolling
  window, and the controls not holding right now. \\
\texttt{direction} & a strategic view & How much of the estate meets the standard its own
  divisions set, and which way that is moving. \\
\texttt{estate} & either & What the estate is made of. Absent values are named
  (\texttt{untiered}, \texttt{unrated}), never dropped. \\
\bottomrule
\end{longtable}

\subsection{Authentication}

A shared token in \texttt{X-ERM-Token}, checked in constant time by the same guard the
ingest surface uses (\texttt{src/web/\_token\_auth.py}). Like ingest's, this is a Phase 0
stub: a real API-key or mTLS scheme replaces it before an external party is onboarded.

The ERM token is deliberately \emph{separate} from the ingest token. The two surfaces have
opposite privileges --- one writes the governance record, one may only read it --- and a
single credential for both would collapse that distinction.

\subsection{Two conventions to code against}

\textbf{Fixed machine values beside adopter labels.} Every enumerated field ships both:
\texttt{state} with \texttt{label}, \texttt{reason} with \texttt{reason\_label}. The
vocabulary is configurable by design (\S\ref{sec:style} makes the same point about colour);
the contract is not. \textbf{Code against \texttt{state} and \texttt{reason}; display
\texttt{label} and \texttt{reason\_label}.} An institution that renames its risk language
must not break its own ERM.

\textbf{Items come back individually.} The console collapses twenty overdue remediation
steps into three rows, one per owner. That is a presentation decision and it stays in the
console. Each item carries \texttt{group} --- the reason it would collapse under, or
\texttt{null} --- so a consumer can reproduce that rollup or list every record instead.
Rolling up inside the contract would leave no way back to the records.

\subsection{What the payload will not do}

Each of these is a commitment, and each is the opposite of what a consumer would otherwise
assume.

\begin{itemize}
  \item \textbf{Report a score it did not compute.}
    \texttt{direction.share\_meeting\_standard} is a count ratio --- records meeting their
    own division's threshold, over records audited. The completeness audit publishes no
    estate-wide figure, and a weighted average invented here would be a number nobody chose
    sitting on a governance record.
  \item \textbf{Infer a deadline.} \texttt{due} is a date a person set, or \texttt{null}. An
    MRC action is agreed \emph{without} a target date, so it is reported as outstanding and
    never as late.
  \item \textbf{Truncate silently.} \texttt{total}, \texttt{returned} and \texttt{truncated}
    are always present.
  \item \textbf{Emit a link that 404s.} \texttt{href} is \texttt{null} where the console has
    no screen for that kind of thing. A dead link in somebody else's dashboard is found by a
    person who cannot fix it.
  \item \textbf{Hide ``no history''.} \texttt{direction.trend} is \texttt{null} until
    observations exist, rather than absent --- a strategic view that cannot tell ``no
    history yet'' from ``flat'' would read a new install as a stable one.
\end{itemize}

\subsection{The contract itself}

JSON Schemas and validated examples are checked in under \texttt{contracts/erm/v1/},
alongside the ingest contract. That is what the ERM's developers code to; the prose here
explains it, but the schema is the agreement. \texttt{tests/web/test\_erm\_api.py} validates
live responses against those files, so a renamed key fails in this repository's suite rather
than in somebody else's dashboard.

\subsection{The finding history}

Ageing a finding needs to know when it first appeared, and the completeness audit is
computed fresh each run with no memory. So a daily observation records what it found:
one row per finding ever seen, carrying first-seen, last-seen and closed-at
(\texttt{0012\_attention\_history.sql}).

Two things follow for an operator. It runs \textbf{on a schedule, from one process} --- the
gunicorn master, beside the backup schedule --- and never on a request, because the console
and this API are both reads. And the two tables it writes are the \emph{only} state in the
platform that is not replayed from the event log: nothing in the log records what the audit
found on a Tuesday, because nobody decided it. \textbf{Back them up.} Losing them loses the
trend line and the ability to age a finding; it loses nothing a committee decided.

\section{Environment variable reference}

Every value below has a shipped default; nothing must be set for the platform to start on a
developer's machine. Real environment variables always beat \texttt{.env}.

\begin{longtable}{@{}p{0.34\textwidth}p{0.24\textwidth}p{0.36\textwidth}@{}}
\toprule
\textbf{Variable} & \textbf{Default} & \textbf{Notes} \\
\midrule
\endhead
\multicolumn{3}{@{}l}{\textbf{Database} --- \texttt{config/database.py}} \\
\midrule
\texttt{DB\_HOST} & \texttt{localhost} & \\
\texttt{DB\_PORT} & \texttt{5433} & The repo's local cluster; a stock server is 5432. \\
\texttt{DB\_NAME} & \texttt{mkm\_modelrisk} & \emph{Protected}: destructive fixtures and
  the verification app refuse it. \\
\texttt{DB\_OWNER\_USER} / \texttt{\_PASSWORD} & \texttt{mkm\_owner} & Owns the schema, runs
  migrations. Keep out of the runtime environment where you can. \\
\texttt{DB\_APP\_USER} / \texttt{\_PASSWORD} & \texttt{mkm\_app} & The runtime role.
  Must be a plain SQL identifier --- the migrations grant to it by name. \\
\texttt{DB\_MAINTENANCE\_NAME} & \texttt{postgres} & Connected to only while creating
  another database. \\
\texttt{TEST\_DB\_BASE} & derived & A dedicated \texttt{*\_test} base name for the suite. \\
\texttt{SCRATCH\_DB} & \texttt{mkm\_modelrisk\_}\newline\texttt{scratch} & Read by
  \texttt{scripts/dev\_app.sh}. \\
\midrule
\multicolumn{3}{@{}l}{\textbf{ERM read API} --- \texttt{config/erm.py}} \\
\midrule
\texttt{ERM\_TOKEN\_HEADER} & \texttt{X-ERM-Token} & \\
\texttt{ERM\_DEV\_TOKEN} & \texttt{dev-erm-token} & \textbf{Change it.} Separate from
  the ingest token on purpose --- see \S\ref{sec:erm}. \\
\texttt{ERM\_ATTENTION\_CAP} & \texttt{500} & Items per response; truncation is
  declared, never silent. \\
\texttt{ERM\_CONSOLE\_BASE} & derived & Where a row's \texttt{href} points. Set it when
  the ERM reaches the console on a different hostname than this API. \\
\midrule
\multicolumn{3}{@{}l}{\textbf{Finding history} --- \texttt{config/history.py}} \\
\midrule
\texttt{HISTORY\_ENABLED} & \texttt{1} & Off loses the trend and finding ages; nothing
  else on the dashboard depends on it. \\
\texttt{HISTORY\_INTERVAL} & \texttt{86400} & Seconds. Daily: the finding set moves at
  the speed of a committee. \\
\texttt{HISTORY\_STARTUP\_DELAY} & \texttt{120} & Seconds, so a restart does not add a
  scan to the stampede. \\
\texttt{TREND\_DAYS} & \texttt{365} & How much series the dashboard and ERM carry. \\
\texttt{CLOSED\_FINDING\_}\newline\texttt{RETENTION} & \texttt{730} & Days a closed
  finding stays on the history. \\
\midrule
\multicolumn{3}{@{}l}{\textbf{Web} --- \texttt{config/web.py}} \\
\midrule
\texttt{WEB\_HOST} & \texttt{127.0.0.1} & \textbf{Set to \texttt{0.0.0.0} in a
  container.} \\
\texttt{WEB\_PORT} & \texttt{5015} & 5016 is the verification app by convention. \\
\texttt{WEB\_BIND} & \texttt{\$WEB\_HOST:\$WEB\_PORT} & As gunicorn spells it. \\
\texttt{WEB\_WORKERS} / \texttt{WEB\_THREADS} & \texttt{4} / \texttt{4} & Peak DB
  connections are their product. \\
\texttt{WEB\_WORKER\_CLASS} & \texttt{gthread} & Requests wait on PostgreSQL; threads absorb
  that cheaply. \\
\texttt{WEB\_TIMEOUT} & \texttt{300} & Generous on purpose --- tectonic renders. \\
\texttt{WEB\_GRACEFUL\_TIMEOUT} & \texttt{30} & \\
\texttt{WEB\_ACCESS\_LOG} / \texttt{WEB\_ERROR\_LOG} & \texttt{-} & The standard streams. \\
\texttt{WEB\_LOG\_LEVEL} & \texttt{info} & \\
\texttt{MODELRISK\_BOOTSTRAP\_DONE} & unset & Set by the master after bootstrap; workers
  inherit it. Do not set it yourself. \\
\texttt{REBUILD\_VIEWS\_ON\_STARTUP} & \texttt{1} & \texttt{0} skips the projection replay
  at boot. \\
\midrule
\multicolumn{3}{@{}l}{\textbf{Identity and permissions} --- \texttt{config/identity.py},
  \texttt{config/permissions.py}} \\
\midrule
\texttt{IDENTITY\_HEADER} & \texttt{X-Auth-Request-User} & \\
\texttt{IDENTITY\_NAME\_HEADER} & \texttt{X-Auth-Request-\newline Preferred-Username} & \\
\texttt{IDENTITY\_DEV\_FALLBACK} & a developer address & Used when no header arrives. Point
  it at a non-person in production. \\
\texttt{SYSTEM\_ACTOR} & \texttt{system} & The platform's own actor for first-run
  seeding. \\
\texttt{ENFORCE\_KNOWN\_ACTORS} & \texttt{protected} & \texttt{always} for a hardened
  deployment; \texttt{never} is break-glass. \\
\texttt{ENFORCE\_REVIEWER\_\newline INDEPENDENCE} & \texttt{always} & \texttt{never} for an
  adopter mid-migration. \\
\midrule
\multicolumn{3}{@{}l}{\textbf{State and backups} --- \texttt{config/files.py},
  \texttt{config/paths.py}, \texttt{config/backup.py}} \\
\midrule
\texttt{UPLOAD\_DIR} & \texttt{<repo>/uploads} & Documents \emph{and} generated reports.
  Must be a volume; must be backed up. \\
\texttt{PG\_DATA\_DIR} & \texttt{./.pgdata} & Only the repo-local cluster. \\
\texttt{PRODUCER\_ROOT} & sibling clone & Root of a producing repository for ingest
  migrations. \\
\texttt{BACKUP\_ENABLED} & \texttt{1} & \texttt{0} for anything that is not the real
  deployment. \\
\texttt{BACKUP\_INTERVAL\_SECONDS} & \texttt{3600} & \\
\texttt{BACKUP\_STARTUP\_DELAY\_SECONDS} & \texttt{20} & So migrations and the rebuild have
  settled first. \\
\texttt{BACKUP\_TIMEOUT\_SECONDS} & \texttt{300} & An unfinished dump is abandoned; the next
  tick retries. \\
\bottomrule
\end{longtable}

\section{Verification checklist}

Run these before calling an integration done. Each one answers a question somebody will ask
later.

\begin{longtable}{@{}p{0.46\textwidth}p{0.48\textwidth}@{}}
\toprule
\textbf{Command} & \textbf{Answers} \\
\midrule
\endhead
\texttt{python -c "from config.database import\newline DatabaseConfig;
  print(DatabaseConfig.from\_env().name)"} &
  Which database this process actually resolved to. Run it before any ad-hoc write. \\
\texttt{pytest} & Unit + integration, on a database the suite derives for itself. \\
\texttt{coverage run -m pytest \&\& coverage report} & The \(\geq\)99\% gate (R3). \\
\texttt{python -m audit} & The coding-rule gate. Read the final
  \texttt{PASS:} / \texttt{FINDINGS:} line. \\
\texttt{python -m audit --pdf} & That tectonic works and the reports render. \\
\texttt{pytest tests/e2e -m e2e} & The browser suite, against a live server on 5016.
  Run one at a time --- it truncates a shared database. \\
\texttt{curl -i http://host/api/models} & That the app serves and reaches the database. \\
\texttt{curl -i -H 'X-Auth-Request-User: nobody@example.com'\newline
  -X POST http://host/api/...} &
  That an unregistered actor is refused with \texttt{403} --- i.e. the known-actor rule is
  in force. \\
\texttt{curl -i -H 'X-Auth-Request-User: someone'\newline http://proxy/...} &
  That your proxy \emph{strips} a client-supplied identity header rather than forwarding
  it. The single most important check in \S\ref{sec:permissions}. \\
\texttt{curl -i http://host/api/erm/v1/}\newline\texttt{model-risk} &
  That the ERM surface refuses an unauthenticated caller with \texttt{401}. If it answers,
  your model-risk position is readable by anyone who can reach the port. \\
\texttt{curl -i -X POST -H 'X-ERM-Token: ...'}\newline
  \texttt{http://host/api/erm/v1/model-risk} &
  That the ERM surface is read-only --- \texttt{405}. \\
\texttt{curl -s -H 'X-ERM-Token: ...' http://host/\newline api/erm/v1/model-risk
  | jq .generated\_at} &
  That an ERM panel can tell how old the figure it is showing is. \\
\bottomrule
\end{longtable}

\section{Where things live}

\begin{code}
config/           every parameter (R1). database, web, identity, permissions,
                  theme, reporting, files, paths, audit, backup, governance/
src/database/     the ONLY module that talks to the store (R6).
                  eventstore.py, projections.py, _connection.py, _migrate.py,
                  migrations/*.sql
src/command/      dispatch.py -- confirm-to-amend, actor attribution, authorize()
src/identity/     resolve_principal(), Principal, the known-actor rule
src/governance/   the domain: model inventory, MRC, products, chain, risk policy,
                  documents, and the LaTeX document builders
src/web/          Flask app factory and blueprints (api_*.py), _startup.py,
                  _theme.py, _helpers.py, _token_auth.py (the machine-surface
                  guard), _history.py (the daily observation schedule),
                  api_dashboard.py, api_erm/ (the read-only ERM contract)
src/static/       js/ and css/ -- the console (R7). No build step.
src/templates/    index.html
audit/            the coding-rule audit: one module per scan, plus pdf.py
scripts/          pg.sh, start.sh, dev_app.sh, restore_scratch.sh, backup_db.sh,
                  and the one-off data migrations
contracts/        the frozen wire contracts: ingest/ (producers write) and
                  erm/ (an ERM reads). JSON Schemas plus validated examples
docs/             coding_rules.md, user_guide.md, permissions_design.md,
                  mssql-immutability.md, windows-setup.md, design/, and this guide
tests/            unit, integration, governance, web, audit, e2e, journey
\end{code}

\vfill
\begin{center}\small\textcolor{mkmchrome}{
\textcopyright{} 2022--2026 MKM Research Labs. The software is licensed under the MIT
License; see \texttt{LICENSE}.}\end{center}

\end{document}
